\section{A taxonomy of surrogate models for energy system optimization}
\label{sec:taxonomy}

Authors typically choose a surrogate family for one of three reasons:
data efficiency on small training sets, ease of integration into a
specific solver class, or familiarity from a neighbouring domain.
Table~\ref{tab:T1-taxonomy} consolidates the families we observe in
the reviewed literature, the type of output they produce (point or
interval/posterior), their compatibility with the dominant optimizer
classes (gradient-based, convex, mixed-integer), and their typical
failure modes. The subsections below describe each family in turn,
focusing on the energy-system applications that have driven its
adoption.

% BEGIN inlined table table_T1_taxonomy
\begin{table*}[!t]
  \centering
  \footnotesize
  \caption{Taxonomy of surrogate model families used for optimization in
  energy-system modeling. ``Output'' is point ($p$) or interval / posterior
  ($i$). ``Optimizer compat.'' indicates natural embedding in
  gradient-based ($\nabla$), convex ($\circ$) or mixed-integer
  ($\square$) solvers. References are representative, not exhaustive;
  the per-entry assignment lives in the curated paper library.}
  \label{tab:T1-taxonomy}
  \renewcommand{\arraystretch}{1.2}
  \begin{tabularx}{\textwidth}{@{}L{0.16\textwidth} C{0.05\textwidth} C{0.08\textwidth} L{0.21\textwidth} L{0.20\textwidth} L{0.22\textwidth}@{}}
    \toprule
    Family & Output & Opt.\ compat. & Strengths & Typical failure modes & Where it dominates (refs.) \\
    \midrule
    Polynomial RSM / PCE
      & $p$,$i$ & $\nabla$ $\circ$ $\square$
      & Closed-form moments; analytic chance constraints; parsimonious
      & Curse of dimensionality; smoothness assumption; sparse-grid cost
      & Chance-constrained AC OPF, stochastic dispatch
        \cite{Muhlpfordt20194806,Muhlpfordt20174490,Koirala20242284,Xu20212696,Engelmann20186188} \\
    Kriging / Gaussian process
      & $p$,$i$ & $\nabla$
      & Calibrated posterior; Bayesian acquisition; small-sample friendly
      & $O(n^3)$ scaling; stationarity assumption; kernel choice sensitive
      & Bayesian optimization for UC, MES design
        \cite{Hu2020,Hu2021671,Volodina2022,Schulte2020,Roth2019214,Liu20222679} \\
    Radial basis functions
      & $p$    & $\nabla$
      & Smooth interpolant; cheap; metaheuristic-friendly
      & Centre placement is heuristic; weak in extrapolation
      & Surrogate-assisted evolutionary OPF and HRES sizing
        \cite{Liu20223726,Mahapatra2022721,Mahapatra20161921,Hussien20211883} \\
    Random forest / GBDT
      & $p$,$i$ & $\square$
      & Categorical inputs; robust to outliers; quantile output
      & Step responses break gradients; piecewise-constant tails
      & Forecasting-as-input for planning; scenario screening
        \cite{Eseye201991463,Faraji2020157284,Hanifi2024,Madhukumar20228891} \\
    Neural networks (MLP / CNN / RNN)
      & $p$ & $\nabla$ ($\square$)
      & Universal approximation; large-data scaling; transfer learning
      & Black-box; calibration drift; overconfidence on OOD
      & End-to-end OPF, demand and price forecasting
        \cite{Chen20244723,Chen2022,Wu2022231,Xu20223066,Guo20222057,Li202359995} \\
    Input-convex / monotone NN
      & $p$ & $\nabla$ $\circ$
      & Preserves convexity; embeddable in LP/QP outer loops
      & Limited expressiveness; harder to train
      & Cost-curve and conversion-loss surrogates in MES dispatch
        \cite{Liu20242534,Wu2024391,Pang2023,Lin2023} \\
    Piecewise-linear MIP-friendly NN
      & $p$ & $\square$
      & Exact embedding into MILP; supports branch-and-bound
      & Big-M tightness; MILP size grows with width
      & MIP-embedded UC, OPF, and capacity expansion proxies
        \cite{Wu2022231,Liu20223726,Liang20246508,Pang2023} \\
    Graph neural networks
      & $p$    & $\nabla$
      & Topology-aware; transferable across grid sizes
      & Message-passing depth; oversmoothing on radial DH/grid trees
      & Power-flow surrogates, multi-fidelity GNN OPF
        \cite{Park2019,Su2022315,Khayambashi202453,Taghizadeh2024} \\
    Physics-informed NN
      & $p$,$i$ & $\nabla$
      & Encodes balance / KKT residuals; better OOD behaviour
      & PDE-loss tuning hard; convergence sensitive to weighting
      & PINN-OPF, building-energy and thermal-network emulators
        \cite{Park2019,Petrova2025,Du2025,Chaudhry2024} \\
    Decision-focused / L2O proxies
      & $p$ & $\nabla$
      & Trained for decision quality, not regression error
      & Feasibility restoration needed; dataset is optimizer-synthetic
      & Learning-to-optimize for SCED / SCUC / OPF
        \cite{Wu2022231,Chen20244723,Chen2022,Liu20223726,Wu2024391,Liang20246508} \\
    \bottomrule
  \end{tabularx}
\end{table*}
% END inlined table table_T1_taxonomy

% =====================================================================
% 3.1 Polynomial response surfaces and polynomial chaos expansions
% =====================================================================

\subsection{Polynomial response surfaces and polynomial chaos
expansions}
\label{sec:tax:rsm}

Polynomial response surfaces have the longest track record in
energy-system optimization. Their main attractions are interpretability,
ease of fit on small designs and direct compatibility with
gradient-based and convex solvers. Response-surface methodology
(RSM) appears in PV control parameter tuning
\cite{Hasanien2018168}, in nano-fluid heat-exchanger design and other
thermal-system optimization problems
\cite{Sami20181416,Ahmad202451,Sun2020}, in chance-constrained
AC optimal power flow as an iterative response-surface refinement
\cite{Xu20212696}, and in green-hydrogen and process-energy
applications \cite{Ahmad202451,Zulfiqar20211026}. The key practical
limitation is the curse of dimensionality: classical low-order
polynomial fits become unreliable above ten to fifteen design
variables \cite{Torun20192128,Meng20191219,Hasanien2018168}.

Polynomial chaos expansions (PCE) generalise RSM to a probabilistic
setting by representing the output as an orthogonal-polynomial series
in random inputs. PCE is the dominant surrogate for chance-constrained
and uncertainty-quantified power-system problems. M\"uhlpfordt et
al.~\cite{Muhlpfordt20194806} pioneered the chance-constrained AC OPF
using PCE; recent work extends the formulation to general PCE-based
hosting-capacity calculations \cite{Koirala20242284}, to data-driven
distribution planning \cite{Fu20202904,Fu2022}, and to efficient
uncertainty quantification in stochastic economic dispatch
\cite{Safta20172535}. PCE delivers analytic moments of the output,
which is exactly what chance-constraint reformulations and
distributionally-robust formulations need; but it scales poorly with
the number of stochastic dimensions and with output non-smoothness
\cite{Muhlpfordt20194806,Safta20172535,Koirala20242284,Xu20212696}.

% =====================================================================
% 3.2 Gaussian process and kriging emulators
% =====================================================================

\subsection{Gaussian process and kriging emulators}
\label{sec:tax:gp}

Gaussian process (GP) regression and its geostatistical variant
kriging are the canonical surrogates for sample-efficient global
optimization and Bayesian uncertainty quantification. They produce a
posterior over outputs at every input, which is what acquisition
functions in surrogate-assisted evolutionary search and Bayesian
optimization use to balance exploitation and exploration.

In dispatch and operations, Nikolaidis et
al.~\cite{Nikolaidis2021} used GP-based Bayesian optimization for
data-driven unit commitment, demonstrating order-of-magnitude
reductions in the number of MILP evaluations needed to converge.
Hu et al.~\cite{Hu2020,Hu2021671} use GP emulators to propagate
uncertainty in stochastic economic dispatch, with a Bayesian-network
extension in \cite{Volodina2022}. Wang et al.~\cite{Wang2023} embed
GP regression in an LSTM-based hybrid for wind-power forecasting and
downstream dispatch, and a comparable hybrid is used by Zhang et
al.~\cite{Zhang2019270}. In Bayesian power-network planning
\cite{Lawson201647} GPs serve as the core uncertainty-aware emulator
for transmission expansion under uncertainty.

In design and expansion problems, kriging-assisted surrogate
evolutionary computation has been used to solve large optimal power
flow instances \cite{Deng2020831} and to optimise battery thermal
management with Gaussian-process-based screening
\cite{Li2021,Babaeiyazdi2021,Kong2021}. The GP family is by far the
most common choice for surrogate-assisted multi-objective and
many-objective design search in microgrids, building retrofits and
energy hubs \cite{Wang2017,Schulte2020,Hussien20211883,Thrampoulidis2021,Roth2019214}.
Cost grows cubically with the training-set size, which is why kernel
sparsification, local GPs and inducing-point variants appear when
data scales beyond a few thousand simulator runs
\cite{Hu2020,Volodina2022,Liu20222679}.

% =====================================================================
% 3.3 Radial-basis-function networks and kernel regressors
% =====================================================================

\subsection{Radial-basis-function networks and kernel regressors}
\label{sec:tax:rbf}

Radial-basis-function (RBF) networks share the same kernel intuition
as GPs but trade the posterior for cheaper deterministic fits. They
are popular in power-system stability classification and in
contingency screening, where the surrogate is queried inside an outer
sampling loop. Urgun et al.~\cite{Urgun20202791} use multilabel
RBF classification for importance-sampled composite-power-system
reliability, and Karamichailidou et al.~\cite{Karamichailidou20212137}
use RBF networks to model wind-turbine power curves for downstream
dispatch. RBF networks also dominate the photovoltaic-power forecast
component in optimization-coupled hybrid pipelines
\cite{Yang2020415,Tian2020}, and they appear as the regression
back-end in extreme-learning-machine load forecasting
\cite{Zeng2017175,Sideratos2020,Madhukumar20228891,Rafi202132436}.
In building-energy applications, kernel regressors have driven the
data-driven heating and cooling load surrogates that feed
district-energy optimisation
\cite{Koschwitz2018134,Hirvonen2017323,Reynolds2017704,Jia2022553}.
Strengths are speed and small-data robustness; the price is poor
extrapolation outside the training cloud, which is why authors who
use RBFs in optimization loops typically combine them with a
trust-region or active-sampling layer
\cite{Deng2020831,Tian2020,Han2018622,Pang2023}.

% =====================================================================
% 3.4 Tree ensembles
% =====================================================================

\subsection{Tree ensembles: random forests and gradient boosting}
\label{sec:tax:trees}

Tree ensembles such as random forests, gradient-boosted trees and
XGBoost are the everyday workhorse of energy-data regression and
classification. Their advantages are tabular data robustness, low
hyperparameter sensitivity, and built-in feature-importance
diagnostics. They appear as predictive components in optimization
pipelines for microgrid scheduling
\cite{Eseye201991463,Faraji2020157284,Wang20192635,Nematollahi2021},
distribution-system planning under uncertainty
\cite{Fu20202904,Fu2022,Jahangiri20244849,Jahangiri2023}, and
short-term wind, solar and load forecasting that drives downstream
dispatch \cite{Li2020,Hanifi2024,Alkesaiberi2022,Madhukumar20228891}.
Tree models do not have natural smooth gradients, which complicates
embedding them in MILP or NLP formulations directly; the practical
workaround is either to use the tree only outside the optimizer
(prediction stage) or to recompile it as a piecewise-linear
big-M MILP
\cite{Ren2024,Lin2023,Pang2023,Wang20211767}.

Tree ensembles also feature in surrogate-assisted multi-objective
optimization where they classify candidate decisions as feasible or
infeasible before more expensive simulator calls
\cite{Han2018622,Schulte2020,Pang2023,Xiao2018}.

% =====================================================================
% 3.5 Neural networks
% =====================================================================

\subsection{Neural surrogates: MLP, recurrent and graph
architectures}
\label{sec:tax:nn}

Neural networks dominate the recent literature, both because of their
flexibility and because of the rise of physics-informed and
graph-aware variants that fit naturally on power and heat networks.

Multi-layer perceptrons (MLPs) are the default choice when the
surrogate has fixed-dimensional input and output and has to be
queried millions of times during an outer optimization loop. MLPs are
used as cost or feasibility proxies in security-constrained
economic dispatch
\cite{Chen2022,Chen20244723,Wu2022231,Xu20223066,Guo20222057,Li202359995},
as the response-surface back-end of metamodel-based design search
\cite{Roth2019214,Hirvonen2017323,Jia2022553,Li2023,Wang2020202933},
and as the surrogate for building retrofit decisions
\cite{Thrampoulidis2021,Reynolds2017704,Koschwitz2018134}.

Recurrent and convolutional architectures (LSTM, GRU, CNN, hybrid
CNN-LSTM) appear when the input is itself a time series and the
surrogate has to learn temporal dependencies. They are widely used
to feed dispatch and storage scheduling
\cite{Rafi202132436,Sideratos2020,Li2019,Tian2020,Yang2020415,Wang2023,Shi2024,Lu20232989,Hanifi2024}.
The downside in optimization use is that recurrent surrogates are
hard to embed in MILP and gradient-based solvers without unrolling.

Graph neural networks (GNNs) are increasingly used as surrogates for
network-structured quantities -- power flows, voltage magnitudes,
contingency severity -- because their inductive bias matches the
topology of the underlying grid
\cite{Park2019,Su2022315,Liu20223726,Chen2023657}. Where the
underlying physics is well understood, physics-informed neural
networks (PINNs) build a residual constraint into the loss, which
sharpens calibration and reduces extrapolation error
\cite{Park2019,Su2022315,Liu20222679,Chen2023657}.

% =====================================================================
% 3.6 Constraint-aware and structure-preserving surrogates
% =====================================================================

\subsection{Constraint-aware and structure-preserving surrogates}
\label{sec:tax:structured}

A central practical issue with neural and tree surrogates is that
the optimization layer has to trust their outputs even on inputs
they have not been trained on. Several recent threads address this
by enforcing structural properties on the surrogate itself.
Convex-mapping formulations
\cite{Roth2019214,Wang2020202933,Hirvonen2017323}, monotone or
input-convex neural networks for energy applications
\cite{Liu20223726,Chen2023657,Chen20244723}, and
piecewise-linear neural networks that can be re-expressed as MILPs
\cite{Wu2022231,Chen2022,Xu20223066,Lin2023} all fall under this
heading. Surrogates that explicitly predict the active-set or the
optimal solution under feasibility guarantees -- end-to-end feasible
optimization proxies \cite{Chen20244723} and learn-to-solve methods
with hard linear constraints \cite{Li202359995} -- close the loop:
the surrogate is no longer just a cheap function evaluator but a
near-feasible decision proposer that the optimizer can validate or
fine-tune. This sub-family is currently the most active part of the
literature and is essential for surrogate-in-MILP use cases that
must respect operational constraints.

% =====================================================================
% 3.7 Hybrid and physics-informed surrogates
% =====================================================================

\subsection{Hybrid and physics-informed surrogates}
\label{sec:tax:pinn}

Several recent contributions combine a cheap physical baseline with
a learned residual on top, which inherits the
extrapolation behaviour of the physical model and the data-fitting
power of the learner
\cite{Park2019,Liu20223726,Chen2023657,Liu20222679,Wang2021,Ren2024,Jalving2023,Zheng20231907}.
Hybrid surrogates are particularly common in district heating and
multi-energy applications, where partial physical knowledge of pipe
dynamics, heat-pump performance maps and storage charging curves is
much cheaper to encode than to learn from scratch
\cite{Schulte2016104,Petrova2025,Gao2025,Reynolds2017704,Jia2022553,Hirvonen2017323,Zhou2024}.
The same idea underpins surrogate-assisted multi-objective
optimization in ship and aircraft energy systems
\cite{Mylonopoulos202332697,Zhang2022,Zhu2024318}.

% =====================================================================
% 3.8 Decision-focused / L2O proxies
% =====================================================================

\subsection{Decision-focused and ``learning to optimize'' surrogates}
\label{sec:tax:l2o}

A distinct paradigm replaces the simulator-of-the-system with the
optimizer-of-the-system: a surrogate is trained to predict the
optimal solution (or a feasible near-optimal one) given a problem
instance. This is the ``learning to optimize'' (L2O) family.
Wu et al.~\cite{Wu2022231} train a deep neural network to solve
security-constrained unit commitment with uncertain wind power.
Chen et al.~\cite{Chen2022,Chen20244723} learn optimization proxies
for large-scale security-constrained economic dispatch with feasibility
guarantees. Guo et al.~\cite{Guo20222057} learn optimal power
allocation in islanded microgrids. Xu et al.~\cite{Xu20223066}
train an ensemble deep network for non-convex economic dispatch.
Li et al.~\cite{Li202359995} learn to solve optimization problems
with hard linear constraints. Liang et al.~\cite{Liang20246508}
combine the L2O idea with statistically feasible robust optimization
for joint chance-constrained unit commitment. Recent reviews track
the explosive growth of this sub-family in power-system
applications \cite{Khaloie2025,Lim2025,Ruan2021221}.
The promise of L2O is the largest speed-up of any surrogate
pattern (orders of magnitude over re-solving the MILP from scratch)
because the surrogate replaces the entire optimizer; the price is
that feasibility and optimality have to be enforced separately,
and that the surrogate is validated under a much tighter standard
than a regression model would be
(Section~\ref{sec:validation}).
